% Education
% =================================================
\DeclareTranslation{french}{section-education}{ÉDUCATION}


% ENSTA
\DeclareTranslation{french}{ensta-title}{Master en Ingénierie Robotique}
\DeclareTranslation{french}{ensta-site}{École Nationale Supérieure de Techniques Avancées (ENSTA), Institut Polytechnique de Paris (IPP)}
\DeclareTranslation{french}{ensta-location}{Palaiseau, Île-de-France, France}
\DeclareTranslation{french}{ensta-duration}{juil. 2022 -- oct. 2025}
\DeclareTranslation{french}{ensta-description}{%
GPA : 3,71/4,00. Programme de double diplôme de trois ans avec l'UNICAMP.
}

% UNICAMP
\DeclareTranslation{french}{unicamp-title}{Licence en Ingénierie de Contrôle et d'Automatisation}
\DeclareTranslation{french}{unicamp-site}{Faculdade de Engenharia Mecânica (FEM), Universidade Estadual de Campinas (UNICAMP)}
\DeclareTranslation{french}{unicamp-location}{Campinas, São Paulo, Brésil}
\DeclareTranslation{french}{unicamp-duration}{fév. 2019 -- août 2026}
\DeclareTranslation{french}{unicamp-description}{%
GPA : 0,84/1,00. Classé 10e sur 59 étudiants et vainquer du Laurier Académique pour mes performances académiques.
}



% Experiences
% =================================================
\DeclareTranslation{french}{section-experiences}{EXPÉRIENCE}


% Ansys
\DeclareTranslation{french}{ansys-title}{Quality Assurance Engineering Intern}
\DeclareTranslation{french}{ansys-site}{ANSYS, Synopsys} % codespell:ignore
\DeclareTranslation{french}{ansys-location}{La Farlède, Provence-Alpes-Côte d'Azur, France}
\DeclareTranslation{french}{ansys-duration}{mai 2025 -- oct. 2025}
\DeclareTranslation{french}{ansys-description}{%
    \begin{itemize}[nosep, leftmargin=*]
        \item \textbf{Automatisation des tests :} Automatisation de 43 tests de régression manuels sous forme de workflows Azure Pipelines à l'aide d'IronPython pour la suite de simulation optique Speos, couvrant des scénarios de capteurs LiDAR et de modélisation de lentilles. Le workflow exécutait les simulations, comparait les sorties et générait des rapports de réussite/échec, éliminant plus de 1 200 minutes de tests manuels par version.
        \item \textbf{Vision par ordinateur :} Développement de zéro d'un utilitaire configurable de comparaison d'images en Python utilisant OpenCV et scikit-image pour les tests de régression, combinant SSIM, détection du flou basée sur le Laplacien, analyse d'histogrammes de couleurs.
    \end{itemize}
}
\DeclareTranslation{french}{ansys-skills}{Azure Pipelines, CI/CD, Git, GitHub Actions, gRPC, IronPython, OpenCV, PyTest, Python, scikit-image}


% SYSNAV
\DeclareTranslation{french}{sysnav-title}{Embedded Firmware and Hardware Designer Intern}
\DeclareTranslation{french}{sysnav-site}{SYSNAV}
\DeclareTranslation{french}{sysnav-location}{Vernon, Normandie, France}
\DeclareTranslation{french}{sysnav-duration}{mai 2024 -- sept. 2024}
\DeclareTranslation{french}{sysnav-description}{%
    \begin{itemize}[nosep, leftmargin=*]
        \item \textbf{Développement de firmware bas niveau :} Développement de zéro d'un driver firmware I2C en C pour un circuit intégré de charge de batterie Texas Instruments au sein d'une base de code STM32H7, avec implémentation de la configuration au niveau des registres, de la gestion des bitfields, de la surveillance d'état et de la gestion des erreurs. Conception de la machine à états finie du driver et validation des protections contre les surintensités, les surchauffes et des transitions d'état de charge conformément à la datasheet du composant.
        \item \textbf{Tests Hardware-in-the-Loop :} Développement d'un framework de tests hardware-in-the-loop en Python avec une STM32H7 Nucleo et 16 relais I2C pour simuler des défauts de connexion du dispositif SYDE. Automatisation du pilotage des relais et de la collecte des logs, réduisant la validation du firmware d'environ 2 jours à 2 heures et couvrant environ 60\% des combinaisons de relais.
    \end{itemize}
}
\DeclareTranslation{french}{sysnav-skills}{Analyseur Logique, C, Git, GitLab, Hardware-in-the-Loop, I2C, Oscilloscope, Python, STM32H7}


% Amazon
\DeclareTranslation{french}{amazon-title}{Retail Vendor Manager Intern}
\DeclareTranslation{french}{amazon-site}{Amazon}
\DeclareTranslation{french}{amazon-location}{Clichy, Île-de-France, France}
\DeclareTranslation{french}{amazon-duration}{sept. 2023 -- mars 2024}
\DeclareTranslation{french}{amazon-description}{%
    \begin{itemize}[nosep, leftmargin=*]
        \item \textbf{Automatisation du traitement des données :} Développement d'un outil VBA de mise en correspondance de mots-clés utilisant des expressions régulières pour classifier des données de licences sur plus de 400 000 entrées, réduisant le temps de traitement de 6 heures à 20 minutes.
    \end{itemize}
}
\DeclareTranslation{french}{amazon-skills}{Analyse de Données, Business Intelligence, Microsoft Excel, Optimisation des processus, SQL, VBA}


% ISIR
\DeclareTranslation{french}{isir-title}{Neuroscience Research Intern}
\DeclareTranslation{french}{isir-site}{ISIR, Sorbonne Université}
\DeclareTranslation{french}{isir-location}{Paris, Île-de-France, France}
\DeclareTranslation{french}{isir-duration}{mai 2023 -- août 2023}
\DeclareTranslation{french}{isir-description}{%
    \begin{itemize}[nosep, leftmargin=*]
        \item \textbf{Simulation sensorimotrice :} Mise en œuvre d'un environnement de simulation Python/Webots basé sur des modèles mathématiques de la cognition spatiale de l'hippocampe, incluant une analyse OpenCV basée sur une caméra pour la cartographie de la navigation.
    \end{itemize}
}
\DeclareTranslation{french}{isir-skills}{Git, LaTeX, Matplotlib, Modélisation Statistique, OpenCV, Python, SciPy, Webots}

% Phoenix
\DeclareTranslation{french}{phoenix-title}{Head of Electronics}
\DeclareTranslation{french}{phoenix-site}{Equipe Phoenix de Robótica da UNICAMP}
\DeclareTranslation{french}{phoenix-location}{Campinas, São Paulo, Brésil}
\DeclareTranslation{french}{phoenix-duration}{juil. 2020 -- juil. 2021}
\DeclareTranslation{french}{phoenix-description}{%
    \begin{itemize}[nosep, leftmargin=*]
        \item \textbf{Leadership :} Direction d'une division électronique de 12 membres, avec développement de formations Altium et en électronique pour les nouveaux membres et coordination du développement, de la fabrication et du dépannage des cartes.
    \end{itemize}
}
\DeclareTranslation{french}{phoenix-skills}{Altium 365, Altium Designer, Débogage matériel, Conception PCB, Fabrication et assemblage.}



% Languages
% =================================================
\DeclareTranslation{french}{section-languages}{LANGUES}


% english
\DeclareTranslation{french}{english-name}{Anglais}
\DeclareTranslation{french}{english-cefr}{C1}
\DeclareTranslation{french}{english-hbar}{83}%
\DeclareTranslation{french}{english-level}{Courant}


% french
\DeclareTranslation{french}{french-name}{Français}
\DeclareTranslation{french}{french-cefr}{C1}
\DeclareTranslation{french}{french-hbar}{83}%
\DeclareTranslation{french}{french-level}{Courant}


% Portuguese
\DeclareTranslation{french}{portuguese-name}{Portugais}
\DeclareTranslation{french}{portuguese-cefr}{C2}
\DeclareTranslation{french}{portuguese-hbar}{100}%
\DeclareTranslation{french}{portuguese-level}{Langue Maternelle}


% Spanish
\DeclareTranslation{french}{spanish-name}{Espagnol}
\DeclareTranslation{french}{spanish-cefr}{B1}
\DeclareTranslation{french}{spanish-hbar}{50}%
\DeclareTranslation{french}{spanish-level}{Intermédiaire}



% Presentation
% =================================================
\DeclareTranslation{french}{presentation}{%
Ingénieur logiciel embarqué avec une expérience pratique du développement de firmware en C pour les microcontrôleurs STM32 et ESP32-S3, notamment de drivers de périphériques, de protocoles de communication et de machines à états finies. Expérience en validation hardware-in-the-loop, débogage matériel, CI/CD et développement robotique. Diplômé d'un master en Ingénierie Robotique de l'ENSTA Paris.
}



% Projects
% =================================================
\DeclareTranslation{french}{section-project}{PROJETS}


% dolphin
\DeclareTranslation{french}{dolphin-title}{Embedded Software and Hardware Developper}
\DeclareTranslation{french}{dolphin-site}{Equipe Paralela -- Dolphin, contrôleur open source pour robot de sumo}
\DeclareTranslation{french}{dolphin-location}{Campinas, São Paulo, Brésil}
\DeclareTranslation{french}{dolphin-duration}{fév. 2026 -- présent}
\DeclareTranslation{french}{dolphin-description}{%
    \begin{itemize}[nosep, leftmargin=*]
        \item \textbf{Firmware :} Développement d'un firmware C pour un contrôleur de robot ESP32-S3 utilisant ESP-IDF, PlatformIO, CMake et une couche FreeRTOS minimale, avec mise en œuvre d'une machine à états finie pilotée par table pour le fonctionnement autonome et radiocommandé. Validation du contrôleur sur un robot de compétition fonctionnel, classé 6e sur 11 robots lors de sa dernière compétition.
        \item \textbf{Matériel :} Conception de la carte PCB principale du contrôleur, du capteur de ligne QRE1113 et des cartes d'encodeur AS5147U sous KiCAD, avec intégration du contrôle moteur PWM, de l'acquisition des capteurs par ADC, du multiplexage commandé par GPIO, des périphériques SPI et des récepteurs RC et capteurs infrarouges pilotés par interruptions.
        \item \textbf{Infrastructure :} Mise en place d'une CI GitHub Actions pour la compilation du firmware et la validation ERC/DRC de KiCAD.
    \end{itemize}
}
\DeclareTranslation{french}{dolphin-skills}{C, ESP32-S3, Git, GitHub Actions, KiCAD, PlatformIO.}
